% Source: https://artofproblemsolving.com/community/c3h1921952_math_problems
% Stable identity: aops-topic-1921952-q05
Write $b=x+1$ for the base and $e=x^2+3x+2=(x+1)(x+2)$ for the exponent. A power $b^e$ equals $1$ only in these cases:
\begin{itemize}
\item $b=1$: then $x=0$, and $1^2=1$. This works.
\item $e=0$ and $b\ne0$: $e=0$ gives $x=-1$ or $x=-2$. For $x=-1$ the base is $0$ and $0^0$ is undefined, so it is rejected. For $x=-2$: $(-1)^0=1$. This works.
\item $b=-1$ and $e$ an even integer: $x=-2$, $e=0$, already found.
\end{itemize}
There are no other cases: if $b>0$ and $b\ne1$, then $b^e=1$ forces $e=0$; if $b<0$, the exponent must be an integer and $|b|^e=1$ forces $|b|=1$ or $e=0$; and $b=0$ gives $0$ or an undefined value. So the solutions are $x=0$ and $x=-2$: there are $\boxed{2}$ real solutions.
